
Mais tout polynôme $P:Q\to\mathbb Q$ est continu dans la topologie induite par $\mathbb R$. Donc pour $k$ assez grand on a
\[
|P_x(n/p^k)-P_x(0)|<1.
\]
Mais par définition $P_x(n/p^k)$ est un entier pour tout $n,k\in\mathbb Z$. Il s'ensuit que $P_x(n/p^k)-P_x(0)=0$ pour $k$ assez grand. Donc le polynôme $P_x$ est constant et on trouve en particulier $P_x(1)=P_x(0)$ i.e. $f(x)=f(0)$.

Les deux cas a) et b) ci-dessus montrent que la restriction de $f$ à $G'$ est constante. On va montrer maintenant que $f$ est constante sur chaque classe de $G$ suivant $G'$. Pour tout $b\in G$ on désigne par $g_b:G'\to\underline{Q}$ l'application définie par
\[
g_b(y)=f(b+y)\qquad(y\in G').
\]
L'application $g_b$ est une application polynôme de $G'$ dans $\mathbb Q$ qui se factorise par $\mathbb Z$, donc constante d'après ce qu'on a vu ci-dessus, d'où l'assertion que $f$ soit constante sur les classes de $G$ suivant $G'$. Il s'ensuit qu'on peut écrire $f=h\circ j$, où $h:G''\to\mathbb Q$ est une application polynôme qui se factorise par $\mathbb Z$. Mais par hypothèse $G''$ est un groupe de torsion de sorte que $h$ est constante. Donc $f$ est constante, C.Q.F.D.

\textbf{Lemme 2.1.4.} \uline{Soit} $A$ \uline{un anneau commutatif unitaire et} $I$ \uline{un idéal nilpotent de} $A$\uline{. Soit} $P$ \uline{le sous-groupe multiplicatif de} $A^*$ \uline{formé des éléments inversibles de} $A$ \uline{de la forme} $x=1+y$ \uline{avec} $y\in I$\uline{, et soit} $i:P\to A$ \uline{l'injection naturelle. Alors l'application composée}
